\section{Conclusion}\label{sec:conclusion}

Day-ahead planning in developing-economy national grids requires synchronised visibility over regional load and system-wide operational pressure, yet forecasting practice often treats demand prediction and stress monitoring as separate pipelines. This study investigated whether a single spatio-temporal framework can jointly forecast next-day regional evening-peak demand and a graph-level Operational Stress Index from shared historical records on the Bangladesh national power system \cite{Hossain2025BangladeshSmartGrid}, under chronological evaluation with classical, recurrent, and graph-convolution baselines and systematic ablation of graph priors, architectural components, and task formulation.

The Correlation-Aware Multi-Task Forecasting Framework implements Parallel-Fusion Spatio-Temporal Graph Transformer (PF-STGT) learning: a seven-day lookback of node-level and national context features over nine divisions, train-estimated correlation-graph conditioning, and parallel regional demand and system-wide stress outputs through shared spatio-temporal encoding and task-specific heads, with a repaired multi-task protocol that stabilises joint learning. The final frozen configuration---correlation-graph PF-STGT with full temporal capacity and multi-task heads---was selected through an evidence-locked programme spanning benchmark comparison, configuration ablation, architecture simplification, error geography, and post-hoc attribution.

Five substantive outcomes emerge. First, graph-transformer multi-task learning delivers operationally coupled day-ahead load and stress forecasts from one checkpoint on a chronologically held-out national timeline. Second, data-driven correlation adjacency outperforms geographical or hybrid graph priors on this dataset, while temporal encoding remains consequential once the correlation prior is adopted. Third, a genuine demand--stress trade-off exists: demand-only training attains a lower megawatt error ceiling, whereas multi-task training enables informative stress outputs at a modest demand offset, positioning the frozen model as the preferred joint optimum. Fourth, regional diagnostics reveal stable error geography---peripheral divisions are comparatively well served while the national load centre dominates absolute residual mass. Fifth, multi-method post-hoc attribution supplies coalition-level, spatial, and temporal transparency with explicit reporting of partial cross-method and dual-path discordance.

Scientifically, the study contributes a reproducible modelling chain on a previously unreported Bangladesh division-level smart-grid case \cite{Musbah2025LoadForecasting,Baur2024XAIReview}, demonstrating that correlation-graph PF-STGT multi-task learning is viable where limitation dynamics and inter-regional load coupling coexist in operational data. Graph priors, task weighting, and architectural simplification require empirical tuning rather than domain intuition; the contribution is bounded to the evaluated national context and frozen internal benchmark suite.

For power-system practice, a single forecasting service may supply paired regional demand trajectories and a national stress scalar for day-ahead review from the same input window. Operators should treat macro summaries as routine briefings, report the principal load centre separately, review limitation and calendar channels highlighted in error and attribution analyses, and present gradient-based and permutation-based explanation views distinctly where methods disagree. These implications are planning diagnostics grounded in offline accuracy and transparency evidence---not demonstrated improvements in grid reliability, economic outcome, or deployment success.

In sum, this work delivers an empirically frozen framework that jointly models regional demand and operational stress on Bangladesh national grid data through correlation-graph PF-STGT learning, supported by a transparent evaluation record from task formulation through architecture selection and post-hoc analysis. Coupled day-ahead load-and-stress forecasting is achievable as a single-model service when graph structure, multi-task calibration, division-level monitoring, and explanation reporting are co-designed with the collected evidence---appropriately qualified as a completed offline evaluation rather than an operational field trial.
